\section{The Continuous-Time Model}\label{sec:model}

\subsection{From an energy bucket to a constant-power circuit (L0, L1)}\label{sec:L0L1}
The simplest reasonable description treats the battery as a bucket of usable energy $E$ drained by the phone's
power demand $P(t)$:
\begin{equation}
  \frac{\dd E}{\dd t} = -\frac{P(t)}{3600}\quad[\mathrm{W\,h\,s^{-1}}],\qquad
  \TTE_0=\frac{E_u}{\bar P}\ \ \text{for constant } \bar P. \label{eq:L0}
\end{equation}
L0 already explains most day-to-day variation: TTE is inversely proportional to mean power. It cannot explain why
the same phone dies early in the cold, why a phone at ``20\,\%'' can shut down, or why gaming time depends on the
room temperature. These effects come from the battery's voltage.

A lithium-ion cell has an open-circuit voltage $\Ocv(\soc)$ that falls with SOC, and a series resistance. The
PMIC's DC--DC converters hold the rail voltages fixed, so the phone draws constant \emph{power}, not constant current
(A1). With $\tilde U=\Ocv-v_1$ and $V=\tilde U-R_0I$, the power balance $P=VI$ gives $R_0I^2-\tilde UI+P=0$. The
physical (high-voltage) root, written without cancellation, is
\begin{equation}
  D=\tilde U^{2}-4R_0P,\qquad I=\frac{2P}{\tilde U+\sqrt D},\qquad V=\frac{\tilde U+\sqrt D}{2}. \label{eq:current}
\end{equation}
This makes the model nonlinear in an important way: as $\Ocv$ falls, the current \emph{rises} to keep the power
constant, which raises the resistive loss and pulls $V$ down further. No solution exists when
$P>\Pmax=\tilde U^2/(4R_0)$; a phone asked for more than $\Pmax$ suffers a voltage collapse. Charge conservation
closes L1:
\begin{equation}
  \frac{\dd\soc}{\dd t}=-\frac{I}{3600\,Q_{\max}},\qquad Q_{\max}=Q_n(1-q_{\mathrm{loss}}). \label{eq:soc}
\end{equation}

\subsection{Component load model}\label{sec:load}
The battery-side power demand is the sum of component powers (A2), each driven by a field of $\bm u(t)$:
\begin{equation}
  P=P_{\mathrm{susp}}+a\,\Delta P_{\mathrm{awake}}+(1-a)\lambda E_{\mathrm{wake}}+P_{\mathrm{disp}}+P_{\mathrm{SoC}}(T_s)
  +P_{\mathrm{net}}+g P_G+c_r\Delta P_{\mathrm{cam}}+c_f\Delta P_{\mathrm{camf}}+au\,\Delta P_{\mathrm{aud}}
  +bt\,\Delta P_{\mathrm{BT}}+P_{\mathrm{sens}}. \label{eq:load}
\end{equation}
Here $a\in\{0,1\}$ means the application processor is awake. The floor terms are the suspend power, the awake
overhead and the energy of background wake-ups while asleep.

\paragraph{Display.} OLED pixel power is linear in linear-light pixel intensity \cite{DZ11}. The content term
therefore scales with the panel area $A$ and the average picture level (APL), and black pixels cost only the fixed
driver overhead $P_{d0}$ \cite{CH13}:
\begin{equation}
  P_{\mathrm{disp}}=s_{\mathrm{scr}}\Big[P_{d0}+p_cA\Big(\frac{b}{0.5}\Big)^{\gamma_b}\frac{m_{\mathrm{APL}}\mathrm{APL}}{\mathrm{APL_{ref}}}
  +\kappa_{RR}\max(f_r-60\,\mathrm{Hz},0)\Big]. \label{eq:display}
\end{equation}
The slider value $b$ maps to luminance as $b^{\gamma_b}$; $m_{\mathrm{APL}}<1$ represents a dark theme and
$\kappa_{RR}$ the cost of high refresh rates \cite{AT20}.

\paragraph{SoC with DVFS.} Dynamic power is $CV^2f$ \cite{LX}. A utilisation-tracking governor sets
$f\propto x_c$ (demanded throughput), and the operating-point voltage is affine in $f$. Hence
\begin{equation}
  P_{\mathrm{SoC,dem}}(x_c)=P_{\mathrm{SoC,max}}\,x_c\,[v_0+(1-v_0)x_c]^2. \label{eq:dvfs}
\end{equation}
Power grows almost cubically: the last 20\,\% of throughput costs 37\,\% of the maximum SoC power.

\paragraph{Radios.} Cellular radios stay in high-power states for a ``tail'' after each transfer
\cite{HU12,NA21}. If bursts arrive as a Poisson process of rate $\lambda$ and restart the timers (A6), the radio
is in the connected-active state exactly when a burst arrived within the inactivity timer $T_{\mathrm{act}}$,
and in the tail state when the last burst arrived between $T_{\mathrm{act}}$ and $T_{\mathrm{tail}}$ ago. The
time-averaged power is then exact:
\begin{equation}
  \Phi=m\Big[(P_{\mathrm{tail}}-P_{\mathrm{idle}})(1-e^{-\lambda T_{\mathrm{tail}}})+(P_{\mathrm{act}}-P_{\mathrm{tail}})(1-e^{-\lambda T_{\mathrm{act}}})
  +\lambda e^{-\lambda T_{\mathrm{tail}}}E_{\mathrm{pro}}+e_b(d_{dl}+\kappa_{ul}d_{ul})\Big], \label{eq:radio}
\end{equation}
where $\lambda e^{-\lambda T_{\mathrm{tail}}}$ is the rate of promotions from idle and $e_b$ the energy per bit.
Real-time media (calls) use a measured connected power $P_{rt}$. Weak signal multiplies cellular traffic energy
by $m=1+\gamma_{\mathrm{sig}}\max(0,(\mathrm{RSRP_{ref}}-\mathrm{RSRP})/10\,\mathrm{dB})$, which reproduces the
$+52$\,\% energy between $-85$ and $-105$\,dBm measured by \cite{DI13}. Wi-Fi, LTE and 5G non-standalone (NSA)
differ only in their parameters (\cref{tab:params}).

\paragraph{Peripherals.} GNSS ($P_G$, \cite{KW19}), rear and front camera pipelines \cite{CX13}, audio and
Bluetooth \cite{CH10} add constant increments while active.

\subsection{Electro-thermal equivalent circuit with throttling (L2)}\label{sec:L2}
\paragraph{Polarisation and OCV.} A first-order Thevenin branch represents diffusion and charge-transfer
overpotentials:
\begin{equation}
  \frac{\dd v_1}{\dd t}=\frac{R_1(T)\,I-v_1}{\tau_1(T)}. \label{eq:rc}
\end{equation}
The OCV uses the Chen--Rinc\'on-Mora form \cite{CRM06} fitted to the LiCoO$_2$/graphite pouch-cell OCV of Ai et
al.\ \cite{AI20} (RMSE 7\,mV), stretched to a modern 4.45\,V charge limit, plus a reversible entropic term:
\begin{equation}
  \Ocv(\soc,T)=-b_0e^{-b_1\soc}+a_0+a_1\soc+a_2\soc^2+a_3\soc^3+\Delta V_{hv}\soc^{p}+(T-T_{\mathrm{ref}})\,k_S\frac{\partial U}{\partial T}(\soc). \label{eq:ocv}
\end{equation}

\paragraph{Temperature-dependent resistance.} Ion transport is thermally activated. With
$a_R(T;E)=\exp[(E/R_g)(1/T-1/T_{\mathrm{ref}})]$,
\begin{equation}
  R_0(\soc,T)=R_{0\infty}\big[1+A_Re^{-\soc/\soc_0}\big]a_R(T;E_{a,R_0})(1+r_g),\quad
  R_1=R_{1,\mathrm{ref}}\,a_R(T;E_{a,R_1})(1+r_g),\quad \tau_1=\tau_{1,\mathrm{ref}}\,a_R(T;E_{a,\tau}), \label{eq:resist}
\end{equation}
where the factor $1+A_Re^{-\soc/\soc_0}$ captures the resistance rise near empty and $r_g$ the growth with aging.
At $-20\,^\circ$C the pack resistance is $3.6\times$ its 25\,$^\circ$C value.

\paragraph{Thermal balance.} All power delivered to the phone becomes heat inside it (A4). The battery adds its
irreversible Joule heat and the reversible entropic heat \cite{BERN85}. A single lumped node (A3) then obeys
\begin{equation}
  C_{\mathrm{th}}\frac{\dd T}{\dd t}=P+I^2R_0+\frac{v_1^2}{R_1}-I\,T\,k_S\frac{\partial U}{\partial T}-\frac{T-\Tamb}{k_{\mathrm{env}}R_{\mathrm{th}}},
  \qquad T_s=\Tamb+\frac{R_{\mathrm{skin}}}{R_{\mathrm{th}}}(T-\Tamb), \label{eq:thermal}
\end{equation}
where $k_{\mathrm{env}}$ represents a case, a pocket or a hand, and $T_s$ is the hot-spot skin temperature
\cite{KANG19,PU23}.

\paragraph{Thermal throttling.} The OS keeps the skin below $T_{\lim}\approx45\,^\circ$C \cite{CHI15}. We model
graded mitigation \cite{AOSP-TH} as a smooth logistic cap on SoC power above a floor $P_{fl}$:
\begin{equation}
  P_{\mathrm{SoC}}=P_{fl}+(P_{\mathrm{SoC,dem}}-P_{fl})\,\sigma\!\Big(\frac{T_{\lim}-\Delta T_{\mathrm{thr}}-T_s}{w_{\mathrm{thr}}}\Big)\ \
  \text{if } P_{\mathrm{SoC,dem}}>P_{fl},\qquad \sigma(y)=\frac1{1+e^{-y}}. \label{eq:throttle}
\end{equation}
The delivered performance ratio is $\rho_{\mathrm{perf}}=x_{\mathrm{eff}}/x_c$, where
$P_{\mathrm{SoC,dem}}(x_{\mathrm{eff}})=P_{\mathrm{SoC}}$. The chain $T\to T_s\to P\to I\to T$ is the feedback
loop of \cref{fig:architecture}.

\paragraph{Bursts and brown-out.} Mean power hides sub-second peaks: radio promotions, CPU boosts and camera starts
draw several watts for milliseconds \cite{HU12,LEE20}. During such a burst $v_1$ is frozen (A5), and the
instantaneous terminal voltage is
\begin{equation}
  V_{pk}=\tfrac12\Big(\tilde U+\sqrt{\tilde U^2-4R_0(P+m_b\Delta P_b)}\Big). \label{eq:burst}
\end{equation}
The phone browns out when $V_{pk}\le V_{\mathrm{sys}}$, the minimum system voltage \cite{QC-SMB}. Because $R_0$
multiplies the burst, a resistance that is harmless at 25\,$^\circ$C becomes decisive in the cold and on aged
packs.

\paragraph{The complete fast system.} For state $\bm x=(\soc,v_1,T)$ and input $\bm u(t)$, equations
\eqref{eq:current}--\eqref{eq:throttle} form the system
$\dot{\bm x}=\bm f(\bm x,\bm u(t),\Tamb;\bm\theta,q_{\mathrm{loss}})$. On the feasible set $\{D>0,\ 0\le\soc\le1\}$
the right-hand side is locally Lipschitz in $\bm x$ and piecewise continuous in $t$. A unique solution therefore
exists until the first event. The time scales range from 20--135\,s (RC) and 7.5\,min (thermal) to hours (SOC),
so the system is only mildly stiff.

\subsection{Aging layer (L3)}\label{sec:L3}
Capacity fade follows square-root laws in time or charge throughput \cite{SCH14,SCM18}. We integrate each loss
$q$ in squared form, $\dd(q^2)/\dd t=k^2(\text{stress})^2$. This is exact for time-varying stress and has no
singularity at $q=0$. Three mechanisms are tracked: calendar loss (SEI growth at rest, accelerated by temperature
and high SOC \cite{KEIL16}), SEI cycle loss (proportional to throughput) and lithium plating (charging at low
temperature or high rate):
\begin{align}
  \frac{\dd q_{\mathrm{cal}}^2}{\dd t_d}&=\big[k_c\,a_{\deg}(T;E_{\mathrm{cal}})\,g(\soc)\big]^2, &
  \frac{\dd q_H^2}{\dd n}&=\big[k_H\,a_{\deg}(T;E_H)\big]^2, &
  \frac{\dd q_L^2}{\dd n_{ch}}&=\big[k_L\,a_R(T;E_L)\,e^{\beta_L(c_{ch}-c_{\mathrm{ref}})}\big]^2, \label{eq:aging}
\end{align}
with $a_{\deg}=1/a_R$, $g(\soc)$ the SOC stress (high when the cell is held near 4.45\,V), $n$ the full
equivalent cycles and $c_{ch}$ the charge C-rate. Then $q_{\mathrm{loss}}=q_{\mathrm{cal}}+q_H+q_L$,
$\mathrm{SOH}=1-q_{\mathrm{loss}}$, and $r_g=\gamma_Rq_{\mathrm{loss}}$. The constants are calibrated so that SOH
reaches 0.80 after 1000 cycles, as stated for current phones \cite{AP-BATT,GOOG}. Aging thus enters a discharge
through two channels: smaller $Q_{\max}$ and larger $R_0,R_1$.

\subsection{SOC outputs and time-to-empty}\label{sec:events}
The fuel gauge reports 0\,\% when the load voltage reaches $\Vcut\approx3.40$\,V \cite{QC-FG,LEE20}, not when
$\soc=0$. The \emph{displayed} SOC is therefore $s=(\soc-\soc_e)/(1-\soc_e)$ with $\Ocv(\soc_e)=\Vcut$. The
problem's glossary defines SOC as remaining \emph{energy} relative to full; we also report the energy fraction
$\varsigma=\int_{\soc_e}^{\soc}\Ocv\,\dd\zeta\big/\!\int_{\soc_e}^{1}\Ocv\,\dd\zeta$. The two differ by at most
2.4 points. TTE is the first time any event function crosses zero:
\begin{equation}
  \TTE=\inf\{t>0:\ \min_i g_i(\bm x(t),\bm u(t))\le0\},\qquad
  \begin{aligned}
  &g_1=\soc,\ \ g_2=V-\Vcut\ \text{(cut-off)},\ \ g_3=V_{pk}-V_{\mathrm{sys}}\ \text{(brown-out)},\\
  &g_4=D\ \text{(infeasible power)},\ \ g_5=T_{\mathrm{shut}}-T\ \text{(thermal shutdown)}.
  \end{aligned}\label{eq:tte}
\end{equation}
The \emph{stranded charge} $s_{\mathrm{off}}=s(\TTE)$ is the displayed SOC at shutdown. It returns once the phone
warms up, because $\soc$ is not reset \cite{LEE20}. Single trajectories are integrated with LSODA
($\mathrm{rtol}=10^{-7}$) segment by segment between input switches, with terminal events. Batches of thousands
of parameter samples use a vectorised RK4 scheme ($h=2$\,s) that agrees with LSODA to within 0.2\,\% of TTE.
